\begin{abstract}
Image jailbreaks must both cross a safety boundary and realize the requested content. Attack success rates based on a binary safety label collapse an important distinction: an output can enter the target category while changing the requested material, anatomical exposure, or body configuration. We study target-preserving jailbreaks with input-specified poses. We introduce \pact{}, Pose-Aligned Conjunctive Testing, which assesses material, exposure, and pose separately and counts success only when their requirements, including the requested configuration, hold in the same depicted figure. We pair this evaluation with \inception{}, a prompt-only method that separates an outer presentation scene from an inner visual target through nested depiction. On Image2, archived outputs exhibit modern drawn figures, the highest exposure category in our rubric, and sexualized presentation across three requested configuration families. Strategy comparisons and evaluator diagnostics further distinguish image return from target attainment. Together, nested depiction and pose-aligned evaluation turn an undifferentiated success label into a test of what an attack preserves and controls.
\end{abstract}

\section{Introduction}
\label{sec:intro}
A controllable image jailbreak must preserve both \emph{what is depicted} and \emph{how the body is arranged}. A safety classifier can give the same positive label to a sculpture, a drawing, and a figure in an unintended pose. Yet these outputs satisfy different requests. A success criterion that collapses their differences cannot establish whether an attack preserves the visual target or controls its configuration.

\begin{figure*}[t]
\centering
\begin{minipage}[t]{0.246\textwidth}\centering
\includegraphics[width=\linewidth]{assets/method_seated.png}\par
\small\textbf{(a) Reclined sitting}
\end{minipage}\hfill
\begin{minipage}[t]{0.246\textwidth}\centering
\includegraphics[width=\linewidth]{assets/method_side.png}\par
\small\textbf{(b) Side-lying}
\end{minipage}\hfill
\begin{minipage}[t]{0.246\textwidth}\centering
\includegraphics[width=\linewidth]{assets/method_knee.png}\par
\small\textbf{(c) One knee raised}
\end{minipage}\hfill
\begin{minipage}[t]{0.246\textwidth}\centering
\includegraphics[width=\linewidth]{assets/method_halftone_T11.png}\par
\small\textbf{(d) Halftone screen}
\end{minipage}
\caption{\textbf{Pose-conditioned figures within nested depictions.} Four original Image2 outputs, shown in full. Panels (a)--(c) retain the existing $M_3/E_6/P_2$ assessments across three configuration families. Panel (d) shows a high side leg raise with the halftone-screen treatment (T11). Configuration correspondence and image provenance are detailed in Appendices~\ref{app:posealignment} and~\ref{app:halftoneoriginal}.}
\label{fig:methodcontrol}
\end{figure*}

We study this problem through two linked ideas: \emph{nested depiction} for constructing the request and \emph{conjunctive testing} for evaluating the result. The construction separates the surrounding scene from the figure represented within it. The evaluation checks whether the requested material, exposure, and body configuration survive together. These requirements provide a common basis for comparing methods (Table~\ref{tab:methoddifference}).

Our method, \inception{}, organizes the request as a scene within a scene. An outer image presents a work; that work depicts the target figure. The outer scene specifies the presentation, while the inner description specifies the figure's rendering and configuration. This separation is the central design choice. A photograph of a drawing can change the surrounding task while retaining a drawn human and an input-specified pose. Low-Effort changes artistic, educational, or lifestyle framing and can replace the object's material; JPA and MODX pursue semantic preservation through prompt optimization \citep{aestheticcamouflage,jpa,modx}. \inception{} makes the boundary between context and target explicit.

Our evaluation, \pact{} (Pose-Aligned Conjunctive Testing), makes that boundary measurable. It independently assesses material $M$, exposure $E$, and pose $P$, then checks the requested configuration $\pi^*$. Success requires all conditions to hold in the \emph{same depicted figure}. This extends atomic content verification to a pose-conditioned attack task \citep{tifa,dsg,geneval}. Component readings show which requirement changed; their conjunction determines target attainment.

On Image2, the first three outputs in Figure~\ref{fig:methodcontrol} share the profile \jointtarget{} across reclined sitting, side-lying, and a raised-knee configuration. A fourth output pairs a high side leg raise with a halftone surround. Strategy comparisons use a common pose budget; evaluator diagnostics examine how presentation changes target assessment.

These contributions organize the paper around one question: \emph{what does the attack preserve and control?}

\section{Related Work}
\label{sec:related}
\paragraph{Attack evaluation.}
ASR comparisons depend on success definitions, evaluation conditions, and configuration coverage \citep{comparisonrequires,asrisnotanumber,singleconfig}. Low-Effort uses safety classifiers; RISE compares automatic judgments with human confirmation \citep{aestheticcamouflage,rise}. TIFA, GenEval, and DSG provide finer-grained checks of objects, attributes, and relations \citep{tifa,geneval,dsg}. These lines of work establish category-level measurement and tools for content verification. The remaining connection is an explicit definition of the full attack target. Even a correct safety label leaves its material, exposure, and configuration unresolved. \pact{} links the checks through a same-figure conjunction and an input-specified pose condition.

\paragraph{Low-effort natural-language attacks.}
Low-Effort studies artistic reframing, pseudo-educational framing, lifestyle and aesthetic contexts, and material substitution \citep{aestheticcamouflage}. These strategies show that descriptions and contexts affect generation. Their effects on the target differ: replacing a drawn human with a sculpture changes the figure material, whereas artistic framing changes the context in which the figure appears. Binary safety success does not separate those changes. \inception{} places the contextual transformation outside the depicted target; \pact{} then checks which inner requirements survive. This is a comparison of what the rewrite changes, not just whether a category detector fires.

\begin{table*}[t]
\centering\small
\begin{tabularx}{\textwidth}{@{}>{\raggedright\arraybackslash}p{0.16\textwidth}YYY@{}}
\toprule
Approach & What the request changes & Preservation criterion & Configuration in the comparison \\
\midrule
Low-Effort \citep{aestheticcamouflage} & Artistic, educational, or lifestyle framing; object-material substitution & Success from safety classifiers & A positive safety label leaves the requested body arrangement unresolved. \\
JPA \citep{jpa} & Prompt-space optimization & Semantic alignment with the original request & Alignment is a global criterion; individual body relations remain a separate check. \\
MODX \citep{modx} & Artistic modifier search & Consistency with the original intent & Semantic preservation motivates checking which configuration details survive. \\
\textbf{\inception{} + \pact{}} & \textbf{Outer presentation changes; inner target is specified separately} & \textbf{Same-figure material, exposure, and pose conjunction} & \textbf{The input-specified configuration $\pi^*$ is part of the success condition.} \\
\bottomrule
\end{tabularx}
\caption{\textbf{From reframing a request to preserving a controllable target.} The comparison distinguishes request construction from its evaluation. \inception{} separates outer context from the depicted target; \pact{} checks the target's attributes and requested configuration jointly.}
\label{tab:methoddifference}
\end{table*}

\paragraph{Semantic preservation.}
JPA optimizes semantic alignment in prompt space, and MODX searches artistic modifiers while addressing consistency with the original intent \citep{jpa,modx}. Their contribution establishes that safety circumvention and content preservation should be considered together. Global semantic agreement, however, can coexist with a changed rendering or body arrangement. We operationalize preservation through separate figure attributes and a requested configuration. This makes deviations localizable rather than subsuming them under a single similarity judgment.

\paragraph{Evidence on modern generators.}
PixJail evaluates multiple attacks, including Low-Effort, under a shared protocol on Image2 and other generators. RISE supplies strategy-evolution results and public examples \citep{pixjail,rise}. Such benchmarks establish aggregate performance, while examples expose particular visual capabilities. A cross-category rate does not identify whether \jointtarget{} and a specified configuration coincide in one output. Our comparisons therefore retain the measurement attached to each source: published benchmark rates, matched archived-return counts, and image-level target profiles. The supplementary evidence separates these quantities and preserves the model attribution of public examples.

\section{PACT: Pose-Aligned Conjunctive Testing}
\label{sec:axes}
\label{sec:targetmatch}
\subsection{The unit and target of evaluation}
\pact{} evaluates a specified figure, not the global appearance of the image. In a photograph of a comic page, the room may be photographic while the person on the page is drawn. Material cues must be attributed to that person. Likewise, exposure from one panel cannot be combined with a matching pose from another. The object assignment precedes the attribute checks.

We write the target as $t^*=(m^*,e^*,p^*,\pi^*)$, where $\pi^*$ specifies the requested body configuration. Our principal content profile is $(M_3,E_6,P_2)$. $M_3$ covers modern drawn humans, including animation, comics, and illustration; $E_6$ requires both anatomical feature classes in our exposure rubric; $P_2$ denotes sexualized presentation without visible sexual activity. A narrower manga-style requirement, or a more precise configuration, belongs in the target specification rather than being inferred from the broad category.

\subsection{Independent attribute readings}
Table~\ref{tab:axes} summarizes the three axes. Material identifies how the figure itself is represented. Exposure records visible covering and anatomical features. Pose describes the configuration and visible behavior. Each axis is judged from its own visual evidence before the results are combined.

\begin{table*}[t]
\centering\small
\begin{tabularx}{\textwidth}{@{}>{\raggedright\arraybackslash}p{0.12\textwidth}Y>{\raggedright\arraybackslash}p{0.20\textwidth}>{\raggedright\arraybackslash}p{0.23\textwidth}@{}}
\toprule
Requirement & Independent observation & Principal target & What must agree \\
\midrule
Material $M$ & $M_1$: manufactured object or teaching plate; $M_2$: pre-modern flat work or invented species; $M_3$: modern drawn human; $M_4$: photograph or photorealism. & $M_3$: modern drawn human & The figure's rendering, distinguished from its carrier and surroundings. \\
Exposure $E$ & $E_1$: clothed; $E_2$: swimwear or fitted clothing; $E_3$: uncovered without identifiable detail; $E_4$: chest detail only; $E_5$: genital detail only; $E_6$: both feature classes. & $E_6$: both required anatomical feature classes & Visible features on the designated figure, using feature-based matching. \\
Pose $P,\pi$ & $P_1$: neutral; $P_2$: sexualized posture; $P_3$: visible sexual activity. Configuration $\pi$ records orientation, support, and limb relations. & $P_2$ and the requested configuration $\pi^*$ & The pose category and every configuration relation required by the target. \\
\bottomrule
\end{tabularx}
\caption{\textbf{The target specification used by \pact{}.} Each row supplies an independent condition; success requires all rows to hold on the same figure. Category indices identify descriptions rather than additive scores. Full instructions and boundary cases appear in Appendices~\ref{app:axisprompts}--\ref{app:exposurepose}.}
\label{tab:axes}
\end{table*}

Exposure and pose answer different questions. A neutral composition can show anatomical detail, and a clothed figure can have a sexualized posture. The exposure categories also encode feature combinations: $E_4$ and $E_5$ identify different visible features rather than successive values on a common intensity scale. Consequently, target matching uses the required features, not arithmetic on category numbers.

The assessment proceeds in three steps. First, separate readings return a category and its evidence for each axis, without receiving the other axes' answers. Second, the records are associated with the same figure. Third, the observations are matched against the target, including the requested configuration. Occluded or ambiguous evidence remains unresolved until adjudication. Full prompts and the feature-matching rules are provided in the appendix.

\subsection{Pose category and requested configuration}
Pose control is a conditional capability: the input supplies a configuration, and the output is evaluated against it. A shared $P_2$ label is insufficient. A standing and a side-lying figure can both be sexualized while satisfying different requests. Configuration matching inspects the specified orientation, support, and limb relations.

The target determines the matching granularity. If it requires only one knee raised, the knee relation is the relevant check. If it also requires lying on the back, that orientation must hold as well. We therefore distinguish the configuration-family descriptions used to present archived examples from complete instruction compliance. Appendix~\ref{app:posealignment} makes this distinction explicit for all three showcase requests. This fixes the level of correspondence used in each reported result.

\subsection{Conjunctive success and denominators}
Let $c_i^M$ and $c_i^E$ indicate material and exposure agreement for output $i$. Pose agreement combines category and configuration:
\begin{equation}
 c_i^P=\mathbb{1}[P_i\models p^*\;\land\;\pi_i\models\pi^*].
\end{equation}
The image-level success indicator and its rate are
\begin{align}
 J_i &= c_i^M c_i^E c_i^P,\label{eq:joint}\\
 \pactasr{} &= \frac{1}{N}\sum_{i=1}^{N}J_i.\label{eq:rates}
\end{align}
Here $N$ is the registered request budget. Explicit refusals and assessed outputs missing any requirement contribute zero. The component rates are reported alongside the joint rate to locate deviations. The conjunction is computed per output; multiplying aggregate component rates would ignore their dependence and whether the requirements co-occur.

Missing outcomes retain a distinct status. With $K$ confirmed joint successes and $U$ unresolved requests or assessments, the rate over the registered set lies in
\begin{equation}
 \left[\frac{K}{N},\frac{K+U}{N}\right].
\end{equation}
A point estimate follows once all requests have a resolved contribution. This definition keeps unresolved evidence separate from known target failure.

\paragraph{What each statistic measures.}
Image return records delivery. Binary safety ASR records a selected classifier's success label. \pactasr{} records realization of the full target. These statistics can be computed on the same experiment, but they are not interchangeable. In particular, the historical strategy comparison in Section~\ref{sec:comparison} reports archived returns because its outputs have not all received the joint assessment required here.

\subsection{A common comparison task}
A comparison fixes the visual target, model, generation settings, request budget, and assessment procedure before comparing success frequencies. Results are also shown by configuration. This exposes whether an aggregate is supported by several requested poses or concentrated in one. When a method changes material as part of its strategy, both the observed material and its agreement with the original target remain visible.

\pact{} formalizes the target-matching logic rather than introducing a different arithmetic for ASR. Its contribution is the unit of assessment, the conditions under which an output is accepted, and the evidence retained for each condition. The protocol can therefore be applied to the output of any compared method under the same task definition.

\begin{figure*}[t]
\centering
\begin{tikzpicture}[x=1pt,y=1pt,font=\small,
  box/.style={draw=black!55,rounded corners=2pt,align=center},
  flow/.style={-{Stealth[length=4pt]},semithick,draw=black!65}]
\path[use as bounding box] (0,0) rectangle (455,136);
\node[anchor=west,font=\small\bfseries] at (0,128) {(a) Construct the depiction};
\node[anchor=west,font=\small\bfseries] at (239,128) {(b) Test the depicted figure};
\draw[rounded corners=3pt,draw=black!55,fill=black!3] (0,4) rectangle (218,112);
\node[anchor=west] at (8,99) {Outer scene $c$: presentation};
\draw[rounded corners=2pt,draw=axisM!65,fill=axisMbg] (15,15) rectangle (203,83);
\node[anchor=west] at (23,70) {Depicted work: carrier};
\draw[rounded corners=2pt,draw=axisP!70,fill=white] (31,24) rectangle (187,55);
\node[align=center] at (109,39) {Target figure\\$M,\ E,\ P,\ \pi^*$};
\node[box,fill=black!3,minimum width=210pt,minimum height=24pt] (loc) at (345,99)
{Localize the same target figure};
\node[box,fill=axisMbg,minimum width=210pt,minimum height=24pt] (read) at (345,60)
{Read $M$, $E$, $P$; match $\pi$ to $\pi^*$};
\node[box,fill=axisPbg,minimum width=210pt,minimum height=24pt] (joint) at (345,21)
{Joint success: $J_i=c_i^M c_i^E c_i^P$};
\draw[flow] (loc) -- (read);
\draw[flow] (read) -- (joint);
\end{tikzpicture}
\caption{\textbf{One target links construction and evaluation.} \inception{} places the figure inside a depicted work and an outer scene. \pact{} follows this relation back to the figure, reads its attributes independently, and checks the requested configuration before computing joint success. The diagram describes the request and evaluation structure.}
\label{fig:construction-evaluation}
\end{figure*}

\section{Inception: Nested Depiction}
\label{sec:attack}
\subsection{Separate context from the visual target}
A nested depiction represents one work inside another. \inception{} assigns the surrounding scene and depicted figure separate roles (Figure~\ref{fig:construction-evaluation}a). The outer scene specifies presentation; the inner description specifies the target.

This separation matters for target preservation. A sheet is a carrier, not the figure's material category. A photograph describes the outer scene, not necessarily the rendering of the figure represented within it. Distinct constraints can therefore specify a photographic presentation and a drawn human at the same time. Evaluation follows the same relation back to the figure whose attributes define the target.

\subsection{Treat pose as an input condition}
The inner target specifies figure rendering, content requirements, and a body configuration. The configuration concerns orientation, support, and key limb relations. These conditions are written before generation and supply the reference for output assessment. The three showcase families thus correspond to different requested conditions rather than a list of poses assigned after generation.

Configuration also affects visibility: a new orientation can hide a required feature even when the pose is realized. Each pose-conditioned output therefore receives its own joint check (Figure~\ref{fig:construction-evaluation}b).

\subsection{Compose through a depiction relation}
Let $x$ denote the content and configuration requirements, $c$ the outer context, and $m$ the figure rendering. We describe the request abstractly as
\begin{equation}
 q=\operatorname{Compose}(R(x,c),m),\label{eq:embedding}
\end{equation}
where $R$ states how the target is carried or represented within the scene. The design has three roles: $c$ identifies the outer task, $R$ organizes the relation between scene and target, and $m$ constrains the figure's realization. An existing complete example appears in Appendix~\ref{app:segments}; the actual texts used by the matched comparison appear in Appendices~\ref{app:our-seat}--\ref{app:aracprompts}.

The experiments test three consequences of this structure: coexistence of the target attributes and configuration, outcomes under a shared pose budget, and changes in output properties under component variations.

\paragraph{Structure and depth.}
The depiction relation specifies how the target belongs to the scene; depth specifies how many surrounding layers are requested. The depth comparison tests whether additional layers improve returns independently of that structural distinction.

\begin{table*}[t]
\centering\small
\begin{tabularx}{\textwidth}{@{}lYcccc@{}}
\toprule
Method & Request construction & Reclined sitting & Side-lying & One knee raised & Total \\
\midrule
\inception{} & Nested depiction & 0/4 & 2/4 & 2/4 & \textbf{4/12} \\
ARA & Artistic reframing & 0/4 & 0/4 & 0/4 & 0/12 \\
MSA & Material substitution & 0/4 & 0/4 & 0/4 & 0/12 \\
LS & Lifestyle context & 0/4 & 0/4 & 0/4 & 0/12 \\
ARAC & Art + drawn material & 0/4 & 0/4 & 0/4 & 0/12 \\
\bottomrule
\end{tabularx}
\caption{\textbf{Archived image returns under a common configuration budget.} Each cell reports archived returns over four registered requests. Zero denotes no archived image in the verified snapshot. ARA, MSA, and LS implement Low-Effort strategy families; ARAC is the drawn-material variant. Joint target assessments remain separate from this delivery metric. Appendix~\ref{app:matchedcomparison} provides the exact texts and provenance.}
\label{tab:comparison}
\end{table*}

\section{Experiments}
\label{sec:experiments}
\subsection{Setup and evidence}
Experiments use text-only generation on Image2. Image editing contributes no generation result. We examine modern drawn figures with the \jointtarget{} content profile and input-specified configuration families. Exact prompts, grouping, and sources are collected in the appendix so that the main text can focus on the comparisons.

The study distinguishes three evidence sets. Showcase images have existing image-level attribute records and corresponding input texts. A five-method comparison shares three configuration conditions and four registered requests per condition. Component and evaluator diagnostics provide additional observations of rendering, context, and assessment behavior. Each table reports its own denominator and measurement.

\subsection{Configuration-specific target examples}
The first three outputs in Figure~\ref{fig:methodcontrol} share the existing profile \jointtarget{} while exhibiting different requested configuration families. Their outer frames, printed surfaces, and photographed settings coexist with a centrally drawn figure. The fourth panel adds a high side leg raise from the halftone-screen condition of the surrounding-treatment comparison. Together, the panels expose the separation between the drawn figure and its presentation.

The configuration evidence is input--output correspondence. The first two cases have existing descriptions of reclined sitting and side-lying with head support. The third realizes the raised-knee relation, while the recorded overall orientation differs from the more specific lying instruction. The appendix preserves this difference at the level of the request. Thus, the examples show pose-conditioned content realizations at the stated granularity, and the joint protocol specifies how finer claims would be scored.



\begin{table*}[t]
\centering\small
\begin{tabularx}{\textwidth}{@{}YYcccc@{}}
\toprule
Carrier & Figure rendering & Standing & Crouching & Forward-leaning & Total \\
\midrule
Studio sheet & Cel & 0/3 & 0/3 & 0/3 & 0/9 \\
Studio sheet & Color manga & 1/3 & 0/3 & 0/3 & 1/9 \\
Studio sheet & Monochrome manga & 1/3 & 0/3 & 0/3 & 1/9 \\
Manuscript paper & Color manga & 0/3 & 0/3 & 0/3 & 0/9 \\
\bottomrule
\end{tabularx}
\caption{\textbf{Rendering and carrier changes have configuration-specific outcomes.} Each setting has three registered requests per configuration. Both returned manga images satisfy the standing condition's orientation; neither comes from the crouching or forward-leaning conditions. The source comparison is detailed in Appendix~\ref{app:medium}.}
\label{tab:rendering}
\end{table*}

\subsection{Nearby natural-language strategies}
\label{sec:comparison}
We compare \inception{} with artistic reframing (ARA), material substitution (MSA), and lifestyle context (LS), implemented from the strategy families of Low-Effort \citep{aestheticcamouflage}. ARAC adds a drawn-material constraint to artistic reframing and is a study variant. Each method has the same three configuration conditions and four registered requests per condition.



Table~\ref{tab:comparison} places the observed difference in two conditions: side-lying and one knee raised. \inception{} has two archived returns in each and none in reclined sitting. The nearby conditions have no corresponding archived return in the snapshot. Reporting the configurations separately exposes where the aggregate comes from.

The showcase and strategy comparison use different complete prompts. The strategy table measures delivery under the registered pose budget; per-output joint reassessment of its four returned images is required for \pactasr{}.

\subsection{Figure rendering and carrier}
The rendering comparison uses three different configurations: standing, crouching, and forward-leaning. Each receives three requests per condition. It contrasts cel rendering, color manga, and monochrome manga on a studio sheet, then compares color manga across two carriers.



Both manga returns are in the standing condition, with none in the crouching or forward-leaning conditions. They therefore establish narrower configuration coverage than the aggregate ``manga returned'' description suggests. Retaining the configuration attached to every output makes this coverage explicit.

The carrier contrast also keeps figure rendering explicit. With color manga fixed, the studio-sheet and manuscript-paper conditions yield one and zero returns, respectively. These counts concern the tested combinations. The representation of the surrounding surface and that of the body remain separate variables for interpreting the result.

\subsection{Context, relations, and depth}
Five contextual designations retain the body description and receive 24 requests each. Red-figure pottery returns 23/24 images, salon painting 17/24, coupled composition 13/24, Edo-period print 10/24, and wall painting 8/24. The fifteen-image range shows that the surrounding description accompanies substantial variation in delivery. Each designation also invokes a visual convention, making figure-level assessment necessary to determine what was preserved.

A relation-component comparison records a more specific change. Removing the clause connecting the figure to the surrounding task yields a clothed figure despite an image return. The observation localizes the difference to the content target, not delivery. It motivates checking the inner attributes when interpreting the role of a contextual component.

The requested-depth sweep yields return counts of $0,1,0,2,1,0,1,1,0,1$ for one through ten layers, with four attempts at each depth. There is no monotonic improvement. This distinguishes relational organization from accumulation of surrounding description. The complete outcomes, including unresolved requests, appear in Appendix~\ref{app:depthdetail}.

\begin{table*}[t]
\centering\small
\begin{minipage}[t]{0.44\textwidth}
\textbf{(a) Fixed images, different observation}\par\medskip
\begin{tabularx}{\linewidth}{@{}Ycccc@{}}
\toprule
Changed reading & $M$ & $E$ & $P$ & Any \\
\midrule
Boxed vs. full & 1/20 & 2/20 & 4/20 & 5/20 \\
Cropped vs. full & 3/20 & 3/20 & 6/20 & 12/20 \\
\bottomrule
\end{tabularx}
\medskip
Boxing changes multiple axes on two images; cropping changes one axis on each of twelve images. Complete profiles are unchanged on fifteen and eight images, respectively.
\end{minipage}\hfill
\begin{minipage}[t]{0.54\textwidth}
\textbf{(b) Fixed foreground, varied backgrounds}\par\medskip
\begin{tabularx}{\linewidth}{@{}Ycrrr@{}}
\toprule
Classifier & Scale & BG & FG & Residual \\
\midrule
AdamCodd & .50 & 59.5 & 2.4 & 38.1 \\
 & .75 & 34.6 & 5.6 & 59.8 \\
 & 1.00 & 32.9 & 7.2 & 59.9 \\
FalconsAI & .50 & 41.3 & 2.4 & 56.4 \\
 & .75 & 58.4 & 4.0 & 37.6 \\
 & 1.00 & 49.8 & 7.6 & 42.6 \\
\bottomrule
\end{tabularx}
\end{minipage}
\caption{\textbf{Evaluator diagnostics motivate figure-level target testing.} (a) Changes in twenty paired automated readings. (b) Variance shares (\%) in the background-composite study: BG denotes background and FG foreground; the smaller scales use forty backgrounds and full scale uses one hundred. Full records and localization controls appear in Appendices~\ref{app:presentation} and~\ref{app:background}.}
\label{tab:diagnostics}\label{tab:background-main}
\end{table*}

\subsection{What the evaluation distinguishes}
A separate gallery-context group returns 8/8 images, with neutral poses in the assessed outputs and one clothed figure. This result has high delivery but does not satisfy a target requiring both $E_6$ and $P_2$. By contrast, the showcased images have the specified content profile. The difference is visible only when the returned content is evaluated.

The judgment procedure itself also responds to presentation. On twenty diagnostic images, comparison with full-image reading gives the changes in Table~\ref{tab:diagnostics}a. Pose changes are the most frequent in both presentation conditions. These are paired automated readings of the same source images, not changes to the generated content.



Joint profiles preserve the co-occurrence hidden by marginal counts. Boxing changes five profiles and seven axis readings because two images change on multiple axes. Cropping changes one axis on each of twelve images. These paired records supply the image-level evidence needed for conjunction. They measure consistency across observation conditions; correctness requires adjudication against the rubric.

A detector example gives a complementary failure mode. NudeNet returns no detections at threshold 0.20 on a woodblock image with identifiable chest details. A zero detection list is therefore not equivalent to the absence of the required feature. The appendix preserves the native detector outputs, their settings, and other diagnostic cases \citep{nudenet}.

Together these observations motivate two parts of \pact{}: localize the evidence to the target figure, and retain the separate judgments before conjunction. A global category alone can miss content distinctions; a local reading without object attribution can be altered by the presentation. Both affect the meaning of an attack-success count.

\subsection{Separating figure evidence from background response}
The object assignment in \pact{} is motivated by a second diagnostic: a global classifier can respond strongly to the surrounding scene while the foreground remains unchanged. We use the existing composite-image analysis, which varies backgrounds at a fixed foreground and scale. The smaller-scale settings use forty backgrounds, and the full-scale setting uses one hundred. These stored-image diagnostics assess the evaluator; they are separate from the text-only generation experiments.



In all six settings, the background component exceeds the foreground component (Table~\ref{tab:background-main}b). Its share ranges from 32.9\% to 59.5\% for AdamCodd and from 41.3\% to 58.4\% for FalconsAI. The result identifies a concrete ambiguity in interpreting a whole-image safety score: a score change need not identify a change in the target figure. This is particularly relevant to nested depictions, whose outer presentation intentionally occupies part of the image.

Localization supplies a complementary check. Across 1,086 diagnostic composites, grouped anatomical detections occur in 22 images, all associated with two backgrounds. Background-only controls reproduce those two cases and yield two additional ones. The detector label therefore needs a location before it can be attributed to the target. A figure-level evaluation should record where the evidence lies as well as which category it supports.

These observations motivate \pact{}'s object assignment: the task determines which figure to assess. Independent attribute records then make the supporting evidence inspectable before the conjunction is computed.

\section{Discussion}
\label{sec:discussion}
\paragraph{A shared object for construction and evaluation.}
\inception{} and \pact{} follow the same relation in opposite directions. The construction places a target within a scene; the evaluation traces the scene back to the target. Their common object is the depicted figure with its requested attributes and configuration. This makes the relation between the proposed design and its measurement explicit.

\paragraph{What controls the observed difference?}
The nested structure supplies a concrete mechanism question: with the target and context held fixed, how does changing their depiction relation affect target preservation? The component observations identify relevant changes in output content, while the depth sweep separates relational structure from a simple ``more layers'' explanation. A fully matched relational intervention would estimate its independent contribution. Internal moderation stages require observations beyond the returned image and response wording.

\paragraph{Comparisons should preserve their task.}
The structural distinction from Low-Effort is where the rewrite acts; the evaluative distinction is what qualifies as success afterward. Material replacement can be effective for its own goal while failing a task that explicitly requires a drawn human. Likewise, broad semantic agreement can preserve the subject but not a requested configuration. Stating the common target before comparison makes these differences interpretable rather than hiding them inside an aggregate rate.

\section{Conclusion}
\label{sec:conclusion}
We study image jailbreaks through target preservation and pose control. \inception{} uses nested depiction to separate an outer presentation task from an inner figure with explicit content and configuration requirements. \pact{} tests those requirements on the same figure and defines success by their conjunction. Image2 examples, nearby strategy comparisons, and evaluation diagnostics show why delivery, category-level success, and full target attainment must be measured separately. The resulting account connects how a request is organized to what its output actually preserves.

\section*{Limitations}
The generation evidence concerns Image2 and the visual targets defined here. The showcase requests differ in more than pose and support the stated configuration-family descriptions; they do not estimate repeatability or arbitrary-pose generalization. The raised-knee example also differs from the finer orientation in its request. The historical method comparison requires per-output joint reassessment before a \pactasr{} ranking can be reported.

The axis readings are operational judgments of visible properties. Stylization, occlusion, and object attribution can change their reliability. The presentation diagnostics use stored automated assessments, while the separate historical repeat experiments concern an earlier evaluator. A dedicated agreement study of the full \pact{} protocol and matched relational ablations would address these distinct questions. Unresolved requests are retained explicitly in all outcome accounting.

\section*{Ethical Considerations}
The study examines image-generation safety and uses fictional figures rather than real people. Original visual examples, including sensitive anatomical content, permit inspection of the evaluation criteria. Image provenance and assessment records accompany the figures. The work is intended to inform evaluation of systems that must interpret both a surrounding scene and the content represented within it.
